\section{Experimental Design for the Extensions}\label{sec:methodology}
\subsection{Questions and Test Populations}
The extended evaluation asks four questions. Can the identification methods
find annotated data without misclassifying code? How often do native and
LLM transformations produce outputs that pass a functional test? How does
wrapping perform when moving from small controlled programs to the original
payload corpus? Finally, what do static scans show once duplicate outputs
and unsuccessful executions are accounted for?

We use three populations with different purposes. The \emph{benign
functional corpus} contains thirteen fixtures, meaning small programs
constructed to exercise specified behavior: seven x64 and six x86 cases
cover arithmetic, branches, loops, local calls, mixed control flow, strings,
and an additional x64 address-acquisition case. Each fixture writes to a
256-byte output buffer. The expected writes are declared before
transformation, and every other byte must retain its initial
\texttt{0xCC} value. This initial value is a \emph{sentinel}: changing it
reveals an unexpected write. A separate C runner returns the complete
buffer for comparison. All thirteen unmodified fixtures pass.

The \emph{classification corpus} contains six further fixtures with
independent reference data ranges: two programs containing multiple strings
(one per architecture), an x64 binary table, two programs containing no data,
and an x64 program whose instruction bytes resemble printable text. This
last case tests whether a classifier mistakes code for a string. Neither
the annotations nor the other classifier's answer is given to the model.
These are six distinct inputs; repetitions measure repeated outcomes on
those inputs rather than enlarging the program population.

The \emph{original payload corpus} contains the same 38 x64 and 64 x86
input files as the conference study. The current harness observes a dialog
for MessageBox, the prescribed connection event for network cases, or
clean process exit for the remaining cases. All 102 originals pass these
checks. These observables are narrower than complete output-buffer
equality: for example, observing a connection does not establish that every
later interaction is correct. We use ``working'' below to mean passing the
specified harness, not proven equivalence for every execution.

\subsection{Versions, Experimental Cells, and Call Accounting}
The code was frozen in three experiment snapshots. B1 (13 November 2026,
13:07 CET) supplies the benign native matrix and x64 LLM tests. B2
(14:05 CET) supplies the x86 LLM tests, x86 wrapping port, and classifier
tests. B3 supplies the full-payload campaign; the hashed rewriting and LLM
modules match B2. Consequently, a comparison between B1 x64 and B2 x86
does not isolate architecture alone.

The benign native matrix uses three runs per input and family. The benign
LLM matrix uses four families, two prompting conditions, and three
single-generation invocations per cell: $7\times4\times2\times3=168$
x64 invocations and $6\times4\times2\times3=144$ x86 invocations.
Rejected responses and transport failures remain in these totals; none is
replaced. Classification uses three LLM invocations and one symbolic run
per annotated fixture.

On the original corpus, wrapping uses three recorded randomization seeds
per input, giving $102\times3=306$ planned variants. It is assessed alone
after data preparation. The LLM matrix has $102\times4\times2=816$ cells,
targeting one accepted output per cell. The output-token ceiling is 65,536.
The saved ledger contains 830 run records: fourteen failed cells were
restarted and failed again. Four run records contain multiple generations;
the other records with known counts contain one. In total, 835 API calls
are recorded, with usage missing for one further run record. Two accepted
outputs required a second generation. Thus this completed campaign is not
a uniform one-generation-per-cell experiment. Cell-level reporting retains
all 816 cells, while cost accounting includes all 830 run records.

The data-identification payload comparison selects nineteen inputs
(thirteen x64, six x86) for which symbolic preparation extracts movable
data. Both methods use the shared deterministic handling and the same
functional harness. This selection evaluates their behavior on those
inputs; it cannot estimate identification accuracy across all 102 inputs.

\subsection{Construction, Functionality, and Identification Measures}
Let $N$ be the number of planned cases, $B$ the number of constructed
outputs, and $W$ the number passing the functional test. We distinguish
\[
  \begin{aligned}
  \text{build rate}&=B/N,\\
  \text{overall working rate}&=W/N,\\
  \text{working rate among builds}&=W/B.
  \end{aligned}
\]
For example, if ten cases yield eight binaries and six pass, these rates
are 80\%, 60\%, and 75\%, respectively. Reporting only the last value
would omit the two cases for which no binary was produced. Retesting a
failed binary does not create a new variant. $U$ denotes the number of
distinct output SHA-256 hashes, a measure of byte-level uniqueness.

For annotated fixtures, exact identification means that all reported data
intervals equal the reference intervals, with no missing or extra region.
We also count missed data bytes and code bytes incorrectly labeled as data.
Writing $P$ for predicted data bytes and $G$ for reference data bytes,
data precision is $|P\cap G|/|P|$ and recall is $|P\cap G|/|G|$.
Inputs without reference data are evaluated as empty-set decisions;
undefined ratios are not replaced by favorable scores. On unannotated
payloads we report interval agreement and final-output identity, not
accuracy. The available comparison matches symbolic extracted ranges
against reported LLM proposals, so it also compares different processing
stages.

\subsection{Representation, Static Scans, and Resource Use}
LLM sequence similarity compares the prepared input listing with an
accepted output using normalized mnemonic-and-operand signatures.
The \texttt{SequenceMatcher} ratio is $S=2M/(u+v)$, where $M$ is the
number of matched elements and $u,v$ are the sequence lengths~\cite{python_difflib}.
A larger value means more sequence overlap, not greater behavioral
correctness. NOP-like instructions are removed in the NOP condition;
therefore $S=1$ can coexist with changed binary bytes. Means include
accepted outputs only. This measure differs from the historical cosine
and entropy measurements in Section~\ref{sect:var_metrics}.

Static scanning has two units of analysis. For a file with $m_i$ malicious
engine responses and $a_i$ total reported responses, $d_i=m_i/a_i$.
Across files, $q$ is the fraction with at least one malicious response.
For example, one malicious response among 75 gives $d_i=1/75$, but the
file counts as detected when computing $q$. The present nonempty reports
contain 75 responses each, including unsupported formats, timeouts, and
failures; those categories are disclosed rather than interpreted as benign
verdicts. An empty report is not an evaluated negative and has no $d_i$.

Scans are deduplicated by hash. A hash shared by multiple transformation
conditions retains membership in each condition, so adding table rows
does not give the size of their union. We also separate passing and failing
LLM outputs before interpreting their scan results. Historical, earlier
pilot, and current scans remain separate cohorts because their inputs,
selection, and collection periods differ.

Time and token usage come from saved records; costs are local tariff
estimates, not invoices. We include failed and restarted runs in campaign
costs and do not treat run duration as pure API latency. The accompanying
evidence freezes the completed logs available on 14 November 2026, including
scan records through 01:17 CET, with source hashes and derived tables.
Earlier benchmark evidence is retained under its original snapshot.
